% M1.tex
\documentclass[11pt,a4paper]{article}
\usepackage[spanish]{babel}
\usepackage[utf8]{inputenc}
\usepackage[T1]{fontenc}
\usepackage{amsmath,amssymb,amsfonts}
\usepackage{siunitx}
\usepackage{graphicx}
\usepackage{booktabs}
\usepackage{hyperref}
\usepackage{tikz}
\usepackage{enumitem}
\usepackage{caption}
\usepackage{mathtools}
\usepackage{float}

\hypersetup{
  colorlinks=true,
  linkcolor=blue,
  citecolor=blue,
  urlcolor=blue,
  pdftitle={Comunicaciones de Datos - Módulo 1},
  pdfauthor={},
  pdfsubject={Apuntes ampliados}
}

\setlist[itemize]{noitemsep, topsep=0pt}
\setlist[enumerate]{noitemsep, topsep=0pt}

% Short commands
\newcommand{\bits}{\mathrm{bits}}
\newcommand{\s}{\mathrm{s}}
\newcommand{\Hz}{\mathrm{Hz}}
\newcommand{\dB}{\mathrm{dB}}
\newcommand{\dBm}{\mathrm{dBm}}
\newcommand{\Watt}{\mathrm{W}}

\begin{document}

\tableofcontents
\bigskip

% -----------------------------
\section{Introducción al modelo de comunicaciones}
El objetivo de estos apuntes es presentar el \emph{modelo de comunicaciones} y los fundamentos físicos y matemáticos que rigen el envío de información entre entidades. En términos generales, comunicar datos consiste en convertir un mensaje en una señal, propagarla por un medio y reconstruir el mensaje en el destino. 

Los bloques funcionales básicos son:
\begin{itemize}
  \item \textbf{Fuente}: genera el mensaje (información) a transmitir.
  \item \textbf{Transmisor}: transforma el mensaje en una señal adecuada para el medio.
  \item \textbf{Sistema de transmisión}: medio físico y nodos intermedios por los que circula la señal.
  \item \textbf{Receptor}: procesa la señal recibida para extraer la información.
  \item \textbf{Destinatario}: actúa como consumidor/interpretador final de la información.
  \item \textbf{Perturbaciones}: ruido e interferencias que distorsionan la señal durante la transmisión.
\end{itemize}

El modelo es \emph{universal} y sirve de referencia para el diseño y análisis de cualquier sistema de comunicación (telefonía, radioenlaces, redes cableadas, etc.).

% -----------------------------
\section{Entidades, mensajes y señales}

\subsection{Entidades}
Una \textbf{entidad} es cualquier elemento capaz de generar o recibir mensajes: computadoras, routers, personas, sensores, etc. En redes normalmente hay bidireccionalidad (full-dúplex o half-dúplex) y, en el modelado, cada sentido puede representarse por una instancia del modelo de comunicaciones.

\subsection{Mensajes}
Un \emph{mensaje} se caracteriza por:
\begin{itemize}
  \item \textbf{Formato cualitativo}: texto, imagen, audio, vídeo o combinaciones (multimedia).
  \item \textbf{Cantidad de información}: medible (ver teoría de la información).
\end{itemize}

\subsection{Señales}
En comunicación de datos, una \emph{señal} es la variación temporal de una magnitud física (eléctrica, óptica, acústica, química, etc.) que representa el mensaje. El envío del mensaje se realiza mediante la propagación de la señal a través de un \emph{medio de transmisión}.

% -----------------------------
\section{Medios de transmisión, tipos de señal y magnitudes}
A continuación se resume la relación típica entre medios, tipo de señal y magnitud física utilizada.

\begin{table}[H]
  \centering
  \caption{Medios de transmisión, tipo de señal y magnitud}
  \begin{tabular}{l l l l}
    \toprule
    Medio & Tipo de señal & Magnitud que varía & Velocidad de propagación (\si{\kilo\meter\per\second}) \\
    \midrule
    Medio metálico (guiado) & Eléctrica & Voltaje / corriente & \SI{230000}{\kilo\meter\per\second} \\
    Fibra óptica (guiado) & Óptica & Intensidad óptica & \SI{214000}{\kilo\meter\per\second} \\
    Espacio libre (no guiado) & EM / óptica & Campo electromagnético & \SI{300000}{\kilo\meter\per\second} \\
    Agua (acústica) & Acústica & Intensidad acústica & \SIrange{1460}{1520}{\meter\per\second} \\
    Cuerpo humano & Química/varios & Concentración molecular & variable \\
    \bottomrule
  \end{tabular}
\end{table}

\paragraph{Notas}
\begin{itemize}
  \item La velocidad de propagación depende del medio; a menudo se expresa como \emph{factor de velocidad} con respecto a la velocidad de la luz en el vacío.
  \item ``No guiado'' se refiere a propagación en espacio libre o medios en los que la señal se dispersa. 
\end{itemize}

% -----------------------------
\section{Tiempo de propagación y tiempo de transmisión}
\subsection{Definiciones}
\begin{itemize}
  \item \textbf{Tiempo de propagación} ($t_p$): tiempo que tarda la señal en recorrer la distancia $d$ por el medio:
  \[
    t_p = \frac{d}{v}
  \]
  donde $v$ es la velocidad de propagación en \si{\meter\per\second}.
  \item \textbf{Tiempo de transmisión} ($t_{tx}$): tiempo necesario para enviar todos los bits del mensaje a la velocidad del enlace $R$ (bits por segundo):
  \[
    t_{tx} = \frac{L}{R}
  \]
  con $L$ la longitud en bits del mensaje.
  \item \textbf{Instante en el que el receptor termina de recibir el mensaje} (si A comienza a enviar en $t=0$): 
  \[
    t_\text{fin} = t_{tx} + t_p.
  \]
\end{itemize}

\subsection{Ejemplos resueltos (ejercicio del PDF)}
Se pide calcular tiempos para un mensaje de \emph{1 MB} entre dos computadores separados por \SI{100}{\meter}.

\medskip
\textbf{Convención usada aquí:} usaremos \emph{1 MB = 10\textsuperscript{6} bytes} = \SI{1e6}{bytes} (formato decimal, frecuente en comunicaciones). Si se quiere usar la definición binaria (MiB = $2^{20}$ bytes) los resultados cambian ligeramente; indico también cómo adaptarlos.

\subsubsection*{Caso (a)}
Dos computadores conectados por cable metálico. Tamaño del mensaje $=1\ \mathrm{MB} = \num{8e6}\ \bits$. Velocidad de transmisión $R=\SI{1e9}{\bits\per\s}$ (\SI{1}{Gbps}). Distancia $d=\SI{100}{\meter}$. Velocidad de propagación en metálico: \SI{230000}{\kilo\meter\per\second} = \SI{230e6}{\meter\per\second}.

\[
t_p=\frac{100}{230\cdot 10^{6}}\ \s \approx 4.3478\times 10^{-7}\ \s \approx 0.435\ \mu\text{s}.
\]
\[
t_{tx}=\frac{8\times 10^{6}}{10^{9}}\ \s = 8\times 10^{-3}\ \s = 8\ \text{ms}.
\]
\[
t_\text{fin}=t_{tx}+t_p\approx 8\ \text{ms} + 0.000435\ \text{ms}\approx 8.000435\ \text{ms}.
\]

Si en cambio se toma $1\ \text{MiB}=2^{20}\ \text{bytes}\approx 1.048576\times 10^{6}\ \text{bytes}$ (bits $\approx 8.388608\times 10^{6}$), el tiempo de transmisión es $\approx 8.3886\ \text{ms}$ y $t_\text{fin}\approx 8.3890\ \text{ms}$.

\subsubsection*{Caso (b)}
Pregunta: ¿en qué instante termina B de recibir el mensaje? Respuesta: $t_\text{fin}$ calculado arriba.

\subsubsection*{Caso (c)}
Mismo escenario pero cable de fibra óptica y velocidad de transmisión $R=\SI{100e9}{\bits\per\s}$ (\SI{100}{Gbps}). Velocidad de propagación en fibra: \SI{214000}{\kilo\meter\per\second} = \SI{214e6}{\meter\per\second}.

\[
t_p=\frac{100}{214\cdot 10^{6}}\ \s \approx 4.6729\times 10^{-7}\ \s \approx 0.467\ \mu\text{s}.
\]
\[
t_{tx}=\frac{8\times 10^{6}}{100\times 10^{9}}\ \s = 8\times 10^{-5}\ \s = 80\ \mu\text{s}.
\]
\[
t_\text{fin}\approx 80\ \mu\text{s} + 0.467\ \mu\text{s} \approx 80.467\ \mu\text{s}.
\]

\subsubsection*{Caso (d)}
Conexión a través de un repetidor instalado en un satélite GEO (altura $\approx$ \SI{36000}{\kilo\meter} sobre el ecuador), velocidad de transmisión $R=\SI{300e6}{\bits\per\s}$ (\SI{300}{Mbps}). Si el enlace A--satélite--B se aproxima por dos tramos verticales de \SI{36000}{\kilo\meter} cada uno, la distancia total recorrida por la señal es $\approx\SI{72000}{\kilo\meter}$.

Velocidad de propagación en espacio libre $\approx \SI{300000}{\kilo\meter\per\second}$.

\[
t_p = \frac{72000}{300000}\ \s = 0.24\ \s = 240\ \text{ms}.
\]
\[
t_{tx}=\frac{8\times 10^{6}}{300\times 10^{6}}\ \s \approx 0.0266667\ \s = 26.6667\ \text{ms}.
\]
\[
t_\text{fin}= t_{tx}+t_p \approx 0.0266667\ \s + 0.24\ \s = 0.2666667\ \s \approx 266.667\ \text{ms}.
\]

\paragraph{Comentarios}
\begin{itemize}
  \item En enlaces de larga distancia (p. ej. satélite GEO) la latencia por propagación domina; incluso con $R$ alto, la propagación introduce retardos grandes.
  \item En enlaces cortos y de alta tasa (fibra 100 Gbps), el tiempo de transmisión domina frente a la propagación en distancias de cientos de metros.
\end{itemize}

% -----------------------------
\section{El modelo de comunicaciones (bloques y ejemplos)}
El sistema de comunicaciones completo incluye además transductores (dispositivos que cambian la naturaleza de la señal, p. ej. micrófono, altavoz, antena, fotodetector).

\subsection{Transductores: tabla ejemplar}
\begin{table}[H]
  \centering
  \caption{Ejemplos de transductores y su conversión}
  \begin{tabular}{l l l}
    \toprule
    Transductor & Entrada & Salida \\
    \midrule
    Micrófono & Acústica & Señal eléctrica \\
    Altavoz & Señal eléctrica & Acústica \\
    Cámara & Óptica (+acústica en multimedia) & Señal eléctrica \\
    Antena & Eléctrica/EM & Campo EM / eléctrica \\
    Lámpara/LED/láser & Eléctrica & Óptica \\
    Fotodetector/fotodiodo & Óptica & Señal eléctrica \\
    Hidrófono & Acústica & Señal eléctrica \\
    \bottomrule
  \end{tabular}
\end{table}

\subsection{Ejemplo: llamada de telefonía fija}
\begin{itemize}
  \item Fuente: persona hablando (señal acústica).
  \item Transductor: micrófono (acústica a eléctrica).
  \item Transmisor/receptor: adaptan la señal a la red pública conmutada (RTC).
  \item Receptor/transductor: altavoz en el destino.
\end{itemize}

\subsection{Ejemplo: radioenlace entre computadores}
\begin{itemize}
  \item Transmisor: modulador + antena.
  \item Medio: espacio libre (propagación electromagnética).
  \item Receptor: antena + demodulador.
\end{itemize}

% -----------------------------
\section{¿Qué estudia la Comunicación de Datos?}
\begin{enumerate}
  \item Fundamentos físicos: medios, transductores, dispositivos, hardware de transmisión.
  \item Fundamentos matemáticos: teoría de la información, teoría de señales, codificación (fuente y canal), modulación.
\end{enumerate}

% -----------------------------
\section{Fundamentos de la Teoría de la Información}
La teoría de la información formaliza la cantidad de información, compresión y límites de transmisión. Las cuatro preguntas clave son:
\begin{enumerate}
  \item ¿Cómo medir la cantidad de información?
  \item ¿Cómo representar la información de forma eficiente (codificación de fuente)?
  \item ¿Cómo proteger la información frente al ruido (codificación de canal)?
  \item ¿Cuál es la máxima tasa que puede transmitir un canal imperfecto (capacidad de canal)?
\end{enumerate}

\subsection{Medida de la información}
Si un evento $x$ tiene probabilidad $p(x)$, la información contenida en la ocurrencia de $x$ se define como:
\[
I(x) = -\log_b p(x),
\]
con base $b$ que determina la unidad (frecuente $b=2$, unidades en \emph{bits}).

\paragraph{Propiedades}
\begin{itemize}
  \item $I(x)\ge 0$.
  \item Si $p(x)=1$ (evento seguro) entonces $I(x)=0$.
  \item Para eventos independientes $x,y$: $I(x,y)=I(x)+I(y)$.
\end{itemize}

\subsection{Entropía de la fuente}
Sea una fuente discreta con alfabeto $\mathcal{X}=\{x_1,\dots,x_M\}$ y probabilidad $p_i=p(x_i)$. La entropía (información media por símbolo) es:
\[
H(\mathcal{X})=-\sum_{i=1}^{M} p_i \log_2 p_i \quad\text{[bits/símbolo]}.
\]
Cumple $0\le H(\mathcal{X}) \le \log_2 M$.

\subsubsection*{Ejemplo (entropía binaria)}
Para una variable binaria con $p$ para el símbolo 1 y $1-p$ para el símbolo 0:
\[
H_2(p) = -p\log_2 p -(1-p)\log_2(1-p).
\]

\subsection{Codificación de fuente}
La codificación de fuente asigna palabras binarias (códigos) a símbolos del alfabeto. Objetivos:
\begin{itemize}
  \item Minimizar la longitud media de código $\bar{L} = \sum_i p_i \ell_i$.
  \item Respetar la decodificabilidad (p. ej. prefijo libre).
\end{itemize}

\paragraph{Eficiencia}
La tasa (bits por símbolo) del código:
\[
R = \frac{\bar{L}}{\log_2 |\mathcal{A}|}
\]
y por el teorema de la codificación sin pérdidas,
\[
H(\mathcal{X}) \le \bar{L}.
\]

\subsubsection{Algoritmos de compresión sin pérdida (resumen)}
\begin{itemize}
  \item Sin memoria: Huffman, Shannon-Fano, codificación aritmética.
  \item Con memoria: Lempel-Ziv (ZIP), BWT, PPM, codificación predictiva.
  \item Formatos comunes: \texttt{ZIP}, \texttt{FLAC} (audio), \texttt{PNG}, \texttt{JPEG-LS} (imágenes sin pérdidas), \texttt{MPEG-4 ALS} (audio).
\end{itemize}

% -----------------------------
\section{Transmisión de la información y canales discretos}
\subsection{Modelo probabilístico}
Se considera la probabilidad absoluta $p(x)$ de que la fuente emita $x$, la probabilidad $p(y)$ de que el receptor reciba $y$, y las probabilidades condicionadas $p(y|x)$ (transición hacia delante) y $p(x|y)$ (transición hacia atrás). Las relaciones básicas incluyen la regla de Bayes y la ley de la probabilidad total.

\subsection{Información mutua}
La información mutua entre entrada $X$ y salida $Y$ (por símbolo) es:
\[
I(X;Y)=\sum_{x}\sum_{y} p(x,y)\log_2\frac{p(x,y)}{p(x)p(y)}.
\]
Otra representación útil:
\[
I(X;Y)=H(X)-H(X|Y)=H(Y)-H(Y|X).
\]

La información mutua media representa cuánta incertidumbre sobre $X$ es resuelta al conocer $Y$.

\subsection{Capacidad de un canal discreto}
La \emph{capacidad} $C$ de un canal discreto es:
\[
C = \max_{p(x)} I(X;Y),
\]
la máxima información mutua por símbolo, optimizando sobre la distribución de entrada.

% -----------------------------
\section{Teoría de la señal (nociones básicas)}
\subsection{Definición}
Una señal $s(t)$ es una función del tiempo que puede ser analógica o digital, periódica o aperiódica.

\subsection{Señal sinusoidal}
Señal senoidal general:
\[
s(t)=A\cos(2\pi f t + \phi),
\]
donde $A$ es la amplitud, $f$ la frecuencia (Hz) y $\phi$ la fase. La potencia media de una sinusoidal con resistencia de carga $R$ es
\[
P_\text{media}=\frac{A^2}{2R}
\]
en términos de amplitud de voltaje pico $A$ (cuando $s(t)$ es voltaje).

\subsection{Potencia de la señal}
Si $s(t)$ es una señal con resistencia equivalente $R$ y corriente/voltaje adecuados, la potencia instantánea $p(t)$ y la potencia media $P$ se definen de forma habitual:
\[
P=\lim_{T\to\infty}\frac{1}{T}\int_0^{T} p(t)\,dt.
\]

\subsection{Decibelio (dB)}
Relación entre dos potencias $P_1$ y $P_2$:
\[
\Delta_{dB} = 10\log_{10}\frac{P_1}{P_2}.
\]
Relación en dBm (referencia 1 mW):
\[
P_{\text{dBm}} = 10\log_{10}\frac{P}{1\ \text{mW}}.
\]

\subsection{Ganancias y atenuaciones en cascada}
En dB las multiplicaciones pasan a sumas. Si hay varios amplificadores con ganancias $G_i$ (factor) y varios elementos con atenuación $A_j$, la potencia de salida total (en dB) es:
\[
P_\text{out,dB} = P_\text{in,dB} + 10\log_{10}\prod_i G_i - 10\log_{10}\prod_j A_j
\]
o, de manera más directa,
\[
P_\text{out,dB} = P_\text{in,dB} + \sum_i G_{i,\dB} - \sum_j \alpha_{j,\dB}.
\]

% -----------------------------
\section{Espectro y respuesta en frecuencia}
\subsection{Relación señal--espectro}
La transformada de Fourier describe la relación entre una señal temporal y su espectro frecuencial:
\[
S(f) = \mathcal{F}\{s(t)\}.
\]
El \emph{ancho de banda} importante de una señal se define como la región de frecuencias donde existe la mayor parte de energía de la señal.

\subsection{Señales periódicas}
Una señal periódica tiene un espectro discreto (con componentes en armónicos de la frecuencia fundamental).

\subsection{Respuesta en frecuencia de dispositivos}
Los dispositivos (filtros, medios, transductores) actúan como ventanas en frecuencia con una respuesta $H(f)$. La salida en frecuencia es $S_\text{out}(f) = S_\text{in}(f)\cdot H(f)$.

% -----------------------------
\section{Transmisión digital y codificación de línea}
\subsection{Transmisión digital}
En redes informáticas la transmisión es digital: las fuentes digitales se conectan normalmente mediante transmisiones digitales (aunque la señal transmitida al medio puede ser analógica —modulada— para transmitir por espacio libre).

\subsection{Codificación de línea}
La codificación de línea adapta una secuencia binaria al canal en banda base. Ejemplos:
\begin{itemize}
  \item NRZ (non-return-to-zero), NRZ polar/unipolar.
  \item Manchester (autorrelojable).
  \item RZ, AMI, HDB3, B8ZS (codificaciones con control de huella DC).
\end{itemize}

\paragraph{Ejemplo: Manchester}
La codificación Manchester invierte el nivel a la mitad del bit: permite sincronización por transición, pero ocupa el doble de ancho de banda nominal.

% -----------------------------
\section{Modulación digital}
\subsection{Concepto}
La modulación digital altera parámetros (amplitud, frecuencia, fase) de una portadora sinusoidal en función de los bits a transmitir. La señal modulada suele ser analógica en banda pasante.

\subsection{Modulaciones binarias básicas}
\begin{itemize}
  \item $2$-ASK / OOK (Amplitude Shift Keying / On-Off Keying).
  \item $2$-FSK (Frequency Shift Keying).
  \item $2$-PSK (Phase Shift Keying) — BPSK.
\end{itemize}

\subsection{Modulaciones multinivel}
\begin{itemize}
  \item QPSK (4-PSK): 2 bits por símbolo.
  \item QAM (Quadrature Amplitude Modulation): combina amplitud y fase; p. ej. 16-QAM, 64-QAM.
\end{itemize}

\subsection{Constelaciones}
Las constelaciones representan en el plano complejo los puntos de señal (amplitud y fase). La distancia entre puntos determina la robustez frente al ruido.

% -----------------------------
\section{Análisis espectral de modulaciones}
La modulación desplaza el espectro de la señal moduladora hacia la banda de la portadora. En sistemas analógicos/paso de banda el ancho de banda necesario depende del tipo de modulación y del filtrado aplicado.

% -----------------------------
\section{Apéndices: fórmulas y tablas de referencia}

\subsection{Conversión entre dBm y mW}
\[
P(\dBm)=10\log_{10} \frac{P(\Watt)}{1\ \text{mW}},\qquad P(\Watt)=10^{(P(\dBm)-30)/10}.
\]

\subsection{Valores típicos}
\begin{itemize}
  \item Potencia de referencia: \SI{0}{\dBm} = \SI{1}{\milli\Watt}.
  \item \SI{30}{\dBm} = \SI{1}{\Watt}.
\end{itemize}

% -----------------------------
\section{Ejercicios propuestos (con indicaciones)}
\begin{enumerate}
  \item Repetir los cálculos del apartado ``tiempos'' suponiendo \emph{1 MiB} en lugar de \emph{1 MB}. Compare los resultados.
  \item Diseñar un código Huffman para la fuente con símbolos $\{a,b,c,d\}$ con probabilidades $\{1/2,1/4,1/8,1/8\}$ y calcular la longitud media del código. ¿Es óptimo?
  \item Para una señal compuesta por varias sinusoidales dadas (amplitudes y frecuencias), calcular su potencia media y representar su espectro.
  \item Calcular la sensibilidad mínima de transmisión en un enlace coaxial con atenuación dada, dadas las ganancias de amplificación y la sensibilidad del receptor (ejercicio similar al del PDF).
\end{enumerate}

% -----------------------------
\section{Notas didácticas y ampliaciones}
\begin{itemize}
  \item Las fórmulas están presentadas en su forma más útil para ingeniería; muchas demostraciones teóricas (por ejemplo, la demostración del teorema de la codificación de fuente y el cálculo de la capacidad de canal) exceden el alcance de estos apuntes, pero son accesibles en textos clásicos (Shannon, Cover \& Thomas).
  \item Cuando se precise más detalle (p. ej. construcción paso a paso de árboles de Huffman, codificación aritmética, o análisis de rendimiento de modulaciones en presencia de ruido AWGN), puedo añadir secciones concretas con ejemplos resueltos y gráficos.
  \item Si deseas, puedo generar también diagramas más detallados con \texttt{tikz} (p. ej. diagramas de bloques o constelaciones) o preparar una versión \texttt{beamer} con las diapositivas replicadas.
\end{itemize}

\bigskip
\noindent\textbf{Fin del documento.}

\end{document}